
\begin{enumerate}[label=\thesubsection.\arabic*, ref=\thesubsection.\theenumi]
    



\item The area of a triangle formed by vertices O, A and B, where 

$\overrightarrow{OA}=$$\hat{i}+2\hat{j}+3\hat{k}$$ $ and $\overrightarrow{OB}=$$-3\hat{i}-2\hat{j}+\hat{k}$$ $ is
\begin{enumerate}
\begin{multicols}{4}
\item $3\sqrt{5}$ sq. units

\item $5\sqrt{5}$ sq. units

\item $6\sqrt{5}$ sq. units

\item $4$ sq. units
\end{multicols}
\end{enumerate}

\solution 

For two given vectors $\overrightarrow{OA}$ and $\overrightarrow{OB}$,

the area of the triangle OAB is given by :
\begin{equation}
Ar(\triangle OAB) = \frac{1}{2}\left|(\overrightarrow{OA} \times \overrightarrow{OB})\right| \tag{1.1}
\end{equation}

so, area of  $\triangle OAB$ is

$$\frac{1}{2}\begin{vmatrix}
\hat{i}&\hat{j} &\hat{k}\\
1&2&3\\
-3&-2&1
\end{vmatrix}$$

$$ \left|4\hat{i}-5\hat{j}+2\hat{k}\right| = \sqrt{4^2 +(-5)^2 +2^2}$$

$$Ar(\triangle OAB)=3\sqrt{5} sq.units$$.

\item

If $cos(sin^{-1}\frac{2}{\sqrt{5}}+cos^{-1}x)=0$, then x is equal to 

\begin{enumerate}
    \begin{multicols}{4}
        
    
\item $\frac{1}{\sqrt{5}}$

\item $-\frac{2}{\sqrt{5}}$

\item $\frac{2}{\sqrt{5}}$ 

\item 1
\end{multicols}
\end{enumerate}
\solution 
\begin{align}
    cos^{-1}x + sin^{-1}x = \frac{\pi}{2}
\end{align}
$cos(\frac{\pi}{2})= 0$ so,
$$cos^{-1}\frac{2}{\sqrt{5}} + sin^{-1}x = \frac{\pi}{2}$$
$$ x= \frac{2}{\sqrt{5}}$$.


\item

The interval in which the function f given by $f(x)=x^{2}e^{-x}$ is strictly increasing is
\begin{enumerate}
    
\begin{multicols}{4}
\item $(-\infty,\infty)$

\item $(-\infty,0)$

\item $(2,\infty)$

\item $(0,2)$
\end{multicols}
\end{enumerate}

\solution 
$$f(x) = x^2e^{-x}$$
$$f^\prime(x) =2xe^{-x} - x^2e^{-x} =0 $$
$2x-x^2 >0$ ,for the $f(x)$ to be strictly increasing

so ,$$x\in (0,2)$$

\item 
The function $f(x)=\frac{x-1}{x(x^{2}-1)}$ is discontinuous at
\begin{enumerate}
    \begin{multicols}{4}
\item exactly one point

\item exactly two points

\item no point

\item exactly three points 
\end{multicols}
\end{enumerate}
\solution 

$f(x) $ domain $\mathbf{R}-\{0,1,-1\}$, so $f(x)$ is discontinuous at $x= 0,1,-1$.

3 points of discontinuity

\item 
The function $f:\mathbf{R}\rightarrow[-1,1]$ defined by $f(x)=cos~x$ is
\begin{enumerate}
    \begin{multicols}{4}
\item both one-one and onto

\item not one-one, but onto

\item one-one, but not onto

\item neither one-one, nor onto
\end{multicols}
\end{enumerate}

\solution 

Range of $cosx$ is [-1,1] , so the function $f(x)$ is surjective or onto
$cosx $ in $\mathbf{R}$ is not one-one.

\item 

The coordinates of the foot of the perpendicular drawn from the point $(2,-3,4)$ on the y-axis is
\begin{enumerate}
    \begin{multicols}{4}
\item $(2,3,4)$

\item $(-2,-3,-4)$

\item $(0,-3,0)$

\item $(2,0,4)$  
\end{multicols}
\end{enumerate}
\solution 

foot of perpendicular of the point  $(x_1,y_1,z_1)$ on the y-axis is $(0,y_1,0) \implies (0,-3,0)$



\item 

The angle between the vectors $i-j$ and $\hat{j}-\hat{k}$ is
\begin{enumerate}
    \begin{multicols}{4}
\item $-\frac{\pi}{3}$

\item 0

\item$\frac{\pi}{3}$

\item$\frac{2\pi}{3}$ 
\end{multicols}
\end{enumerate}

\solution 

angle between two vectors $a_1\hat{i} + b_1\hat{j} +c_1\hat{k} $ and $a_2\hat{i} +b_2\hat{j} +c_2\hat{k}$ is 

\begin{equation}
    cos\theta = \left( \frac{\overrightarrow{u}\cdot \overrightarrow{v}}{\|\mathbf{u}\|\|\mathbf{v}\|}\right) \tag{7.1}
\end{equation}


$cos\theta = \left( \frac{-1}{2}\right)$ , so $\theta =\frac{2\pi}{3}$

\item 

If A is a non-singular square matrix of order 3 such that $A^{2}=3A$ then value of $\vert{A}\vert{}$ is
\begin{enumerate}
    \begin{multicols}{4}
\item -3

\item 3

\item 9

\item 27
\end{multicols}
\end{enumerate}
\solution

$\vert{A^2}\vert{} = \vert{3A}\vert{}$ , $\vert{A}\vert{} \neq 0$

$\vert{A}\vert=27$

\item

If $\vert{}\vec{a}\vert{}=4$ and $-3\le\lambda\le2$, then $\vert{\lambda \overrightarrow{a}}\vert{}$ lies in
\begin{enumerate}
    \begin{multicols}{4}
\item $[0, 12]$

\item $[2, 3]$

\item $[8,12]$

\item $[-12,8]$
\end{multicols}
\end{enumerate}
\solution

$-3\le\lambda\le2$ ,then $0\le\vert{}\lambda\vert{}\le 3$ 

$\vert{}\lambda \overrightarrow{a}\vert{} $ , $4\vert{}\lambda\vert{} \in [0,12]$

so,$$\vert{}\lambda\overrightarrow{a}\vert{} \in [0,12]$$

\item

If the radius of the circle is increasing at the rate of 0.5 cm/s, then the rate of increase of its circumference is   

\solution

$\frac{dr}{dt} =0.5 cm/s$ , circumference $c=2\pi r$

$\frac{dc}{dt}= 2\pi \frac{dr}{dt}$ , $\frac{dc}{dt} = \pi cm/s$

\item 

If $\vert{}\begin{matrix}2x&-9\\ -2&x\end{matrix}\vert{}=\vert{}\begin{matrix}-4&8\\ 1&-2\end{matrix}\vert{}$ then value of x is   

\solution

$2x^2 - 18 =0$ , $x=\pm 3$

\item 

The corner points of the feasible region of an LPP are $(0,0)$ $(0,8),$ $(2,7),$ $(5,4)$ and $(6,0)$ The maximum profit $P=3x+2y$ occurs at the point 

\solution

To find the point at which the maximum value of the objective function $P = 3x + 2y$ occurs, we evaluate $P$ at each of the given corner points[cite: 4]:

\begin{enumerate}
    \item At $(0,0)$:
    $$P = 3(0) + 2(0) = 0$$

    \item At $(0,8)$:
    $$P = 3(0) + 2(8) = 16$$

    \item At $(2,7)$:
    $$P = 3(2) + 2(7) = 6 + 14 = 20$$

    \item At $(5,4)$:
    $$P = 3(5) + 2(4) = 15 + 8 = 23 \quad \text{(Maximum)}$$

    \item At $(6,0)$:
    $$P = 3(6) + 2(0) = 18$$
\end{enumerate}



Comparing all the values, the maximum value of $P$ is $23$, which occurs at the point $(5,4)$



The maximum profit occurs at the point $(5,4)$

\item 

The distance between parallel planes $2x+y-2z-6=0$ and $4x+2y-4z=0$ is units

\solution

$4x+2y-4z= 12$ and $4x+2y-4z=0$ 

distance between parallel planes $ax +by +cz= d_1 $ and $ax+by +cz=d_2$ is $D=\left\vert{}\frac{d_1-d_2}{\vert|u|\vert{}}\right\vert{}$

$D=\left\vert{}\frac{12}{\sqrt{4^2 + 2^2+ 4^2}}\right \vert{}$ , D =2 .

\item 

Evaluate: $\int_{-\frac{\pi}{2}}^{\frac{\pi}{2}}xcos^{2}xdx$

\solution

I =$\int_{-\frac{\pi}{2}}^{\frac{\pi}{2}}xcos^{2}xdx$

I= $\int _{\frac{-\pi}{2}}^{\frac{\pi}{2}}\left(\frac{\pi}{2}-\frac{\pi}{2}-x\right)cos^2 xdx$, I = -I ; I=0.

\item

Find the coordinates of the point where the line $\frac{x-1}{3}=\frac{y+4}{7}=\frac{z+4}{2}$ cuts the xy-plane.   

\solution

z=0 , $$\frac{x-1}{3}= \frac{y+4}{7}=2$$ $$x=7 ,y=10$$

\item 

Find the value of k, so that the function $f(x)= \begin{cases}
    kx^2+5 &ifx\le 1\\ 2 & ifx>1\end{cases} $ is continuous at $x=1$

\solution

$f(1^-)= f(1^+)$ $$k+5 = 2$$ $$k=-3$$

\item 

Find the integrating factor of the differential equation $x\frac{dy}{dx}=2x^{2}+y$

\solution

$x\frac{dy}{dx}-y= 2x^2$ $\implies$ $\frac{dy}{dx}-\frac{y}{x} = 2x $

$$ye^{\int -\frac{1}{x}dx} = \int 2xe^{\int -\frac{1}{x}dx}dx$$

integrating factor  $e^{\int -\frac{1}{x}dx}$


\item

If $y=f(x^{2})$ and $f^{\prime}(x)=e^{\sqrt{x}}$ then find $\frac{dy}{dx}$

\solution

$\frac{dy}{dx}= 2xf^\prime(x^2) \implies \frac{dy}{dx}= 2xe^x$


\item

\textbf{a.}Find a vector $\vec{r}$ equally inclined to the three axes and whose magnitude is $3\sqrt{3}$ units. 

\solution

$\alpha =\beta=\gamma= \cos^{-1}(\pm\frac{1}{3})$

$\overrightarrow{r}=3\sqrt{3}\left(\pm\frac{1}{\sqrt{3}}\hat{i} \pm\frac{1}{\sqrt{3}}\hat{j} \pm\frac{1}{\sqrt{3}}\hat{k}\right)$ 

\begin{center}
    or
\end{center}


\textbf{b.}Find the angle between unit vectors $\vec{a}$ and $\vec{b}$ so that $\sqrt{3}\vec{a}-\vec{b}$ is also a unit vector.   

\solution

$\vert{}\sqrt{3}\overrightarrow{a}-\overrightarrow{b}\vert{}=\sqrt{3\vert{a}\vert{}^2 + \vert{b}\vert{}^2 +2\sqrt{3}\left(\overrightarrow{a}\cdot\overrightarrow{b}\right)}$

$1 = \sqrt{3 + 1 +2\sqrt{3}\left(\overrightarrow{a}\cdot\overrightarrow{b}\right)}$

$\overrightarrow{a}\cdot\overrightarrow{b}= -\frac{\sqrt{3}}{2} = cos\theta \implies \theta = \frac{5\pi}{6}$

\item

If $A=\left[\begin{matrix}-3&2\\ 1&-1\end{matrix}\right]$ and $I=\left[\begin{matrix}1&0\\ 0&1\end{matrix}\right]$ find scalar k so that $A^{2}+I=kA.$

\solution

$\vert{}A-\lambda I\vert{}= 0 \implies \left|\begin{matrix}
    -3-\lambda & 2\\ 1& -1-\lambda\end{matrix} \right|= 0$

    $\lambda^2+ 4\lambda +1=0 \implies A^2 +I= -4A \implies k=-4$


\item

\textbf{a.}If $f(x) =\sqrt{\frac{secx - 1}{secx + 1}}$, find $f^\prime {(\frac{\pi}{3})}$


\solution

$f(x)= \sqrt{\frac{1-cosx}{1+cosx}} \implies f(x)= tan\frac{x}{2}$

$f^\prime (x) = \frac{1}{2}{sec^2\frac{x}{2}} \implies  f^\prime(\frac{\pi}{3}) = \frac{2}{3}$

\begin{center}
    or 
\end{center}


\textbf{b.} Find $f^\prime(x)$ if $f(x)=(tan~x)^{tan~x}$

\solution

$f^\prime(x)= tanx(tanx)^{tanx -1}(sec^2x) + tanx^{tanx}ln(tanx)sec^2x$

$f^\prime(x)= tanx^{tanx}sec^2x\left(1 +ln(tanx)\right)$

\item

Find $\int \frac{tan^3x}{cos^3x}dx$

\solution

cosx=t $ \implies \int \frac{1-t^2}{t^6}dt \implies \int \frac{1}{t^4}dt-\frac{1}{t^6}dt$

$\frac{1}{5cos^5x}-\frac{1}{3cos^3x}$

\item 

Show that the plane $x-5y-2z=1$ contains the line $\frac{x-5}{3}=y=2-z.$ 

\solution

$\frac{x-5}{3}=y=2-z=k \implies x=3k+5 ,y=k,z=2-k$

substituting them in the plane equation gives, \newline $3k+5-5k-2(2-k)-1 =0 \implies 0=0 $, so the line lies in the plane 

\item

A fair dice is thrown two times. Find the probability distribution of the number of sixes. Also determine the mean of the number of sixes.

\solution

X,Y are the outcomes of each die,

$P(X,Y)=\begin{cases}
    \frac{25}{36}& X\neq6,Y\neq6 \\ \frac{5}{18}&X=6,Y\neq6 (or )X \neq 6 ,Y=6 \\ \frac{1}{36} & X=6 ,Y=6\end{cases}  $
    
mean $\mu= E(x) =\Sigma X,P(X)$

$E(x)= (0* \frac{25}{26})+(1*\frac{5}{18})+(2*\frac{1}{36})$

$$E(x)=\frac{1}{3}$$

\item 

Solve the following differential equation \newline $(1+e^{y/x})dy+e^{y/x}(1-\frac{y}{x})dx=0(x\ne0).$

\solution

Let $\frac{y}{x} =t \implies \frac{dy}{dx} =t + x\frac{dt}{dx} $

$$(1+e^t)(t+x\frac{dt}{dx})+ e^t(1-t) =0$$  $$\frac{e^t +t}{e^t+1}dt=-\frac{1}{x}dx$$

$$ln(e^t + t)=ln(\frac{1}{x}) \implies e^t + t= \frac{1}{x} $$

$$e^{\frac{y}{x}} + \frac{y}{x} = \frac{1}{x}$$

\item 

Evaluate : 

$$ \int_0^\frac{\pi}{2}sin2xtan^{-1}(sinx)dx$$

\solution

$I=\int_0^\frac{\pi}{2} sin2xtan^{-1}(sinx)dx=\int_0^\frac{\pi}{2}sin2xtan^{-1}(cosx)dx$

$sinx=t; cosxdx=dt$

$I=2\int_0^1ttan^{-1}tdt \implies  2\left(\left(\frac{t^2}{2}tan^{-1}t\right)_0^1-\int _0^1 \frac{t^2}{2(t^2+1)}dt\right)$

$I =\frac{\pi}{2}-1$

\item 

Prove that $tan^{-1}\frac{1}{4}+tan^{-1}\frac{2}{9}=\frac{1}{2}sin^{-1}(\frac{4}{5})$

\solution

$tan^{-1}\left(\frac{\frac{1}{4} +\frac{2}{9}}{1-\frac{1}{18}}\right)$

$tan^{-1}(\frac{1}{2})=sin^{-1}(\frac{1}{\sqrt{5}})$

$2sin^{-1}(\frac{1}{\sqrt{5}})= sin^{-1}(\frac{4}{5}) \implies sin^{-1}(\frac{1}{\sqrt{5}})=\frac{1}{2}sin^{-1}(\frac{4}{5})$

Hence proved 

\item 

Find the equation of the plane passing through the points $(1,0,-2)$, $(3,-1,0)$ and perpendicular to the plane $2x-y+z=8$ Also find the distance of the plane thus obtained from the origin.  

\solution

$$dr_1 = (2,-1,2); dr_2= (2,-1,1)$$  $$dr=dr_1 \times {dr_2 } \implies (1,2,0)$$

$1(x-1) +2(y) =0 \implies x+2y=1$

\item 

If $tan^{-1}(\frac{y}{x})=log\sqrt{x^{2}+y^{2}}$ , prove that $\frac{dy}{dx}=\frac{x+y}{x-y}$

\solution

$\frac{y}{x}=t \implies \frac{xdy-y}{x^2} =dt$

$tan^{-1}t=\frac{1}{2}(2logx +log(1+t^2)) \implies \frac{1}{1+ t^2}dt= \frac{1}{x}+\frac{tdt}{1+t^2}$

$$\frac{1-t}{1+t^2}dt =\frac{1}{x}$$  $$\frac{x(x-y)}{x^2+y^2}\left(\frac{x\frac{dy}{dx}-y}{x^2}\right)=\frac{1}{x}$$ $$\frac{dy}{dx}= \frac{x+y}{x-y}$$

\begin{center}
    or 
\end{center}
\textbf{question :}

If $y=e^{a~cos^{-1}x},-1<x<1$, then show that $(1-x^{2})\frac{d^{2}y}{dx^{2}}-x\frac{dy}{dx}-a^{2}y=0$

\solution

$\frac{dy}{dx}=-e^{a cos^{-1}x}a(\frac{1}{\sqrt{1-x^2}}) $$\frac{dy}{dx}= \frac{-ay}{\sqrt{1-x^2}} \implies \sqrt{1-x^2}\frac{dy}{dx}= -ay$

$$(1-x^2)(\frac{dy}{dx})^2=a^2y^2$$  $$2\frac{d^2y}{dx^2}(\frac{-ay}{\sqrt{1-x^2}})(1-x^2)-2x\frac{a^2y^2}{1-x^2}= 2a^2y(\frac{-ay}{\sqrt{1-x^2}}) $$

$$(1-x^{2})\frac{d^{2}y}{dx^{2}}-x\frac{dy}{dx}-a^{2}y=0$$

\item 

Amongst all open (from the top) right circular cylindrical boxes of volume $125\pi~cm^{3}$ find the dimensions of the box which has the least surface area.  

\solution

$V=\pi r^2h \implies h= \frac{125}{r^2}$

The total surface area $S$
\begin{align*}
S &= \pi r^2 + 2\pi r h \\
S(r) &= \pi r^2 + 2\pi r \left(\frac{125}{r^2}\right) = \pi r^2 + \frac{250\pi}{r}
\end{align*}

Differentiating with respect to $r$:
\begin{align*}
\frac{dS}{dr} &= 2\pi r - \frac{250\pi}{r^2}
\end{align*}

$$\frac{dS}{dr} = 0$$:
\begin{align*}
2\pi r - \frac{250\pi}{r^2} &= 0 \\
2\pi r^3 &= 250\pi \\
r^3 &= 125 \implies r = 5\text{ cm}
\end{align*}

Second derivative test:
\begin{align*}
\frac{d^2S}{dr^2} &= 2\pi + \frac{500\pi}{r^3}
\end{align*}
At $r = 5$:
\begin{align*}
\left.\frac{d^2S}{dr^2}\right|_{r=5} &= 2\pi + \frac{500\pi}{125} = 2\pi + 4\pi = 6\pi > 0
\end{align*}
the surface area is minimum when $r = 5\text{ cm}$.

Corresponding height:
\begin{align*}
h &= \frac{125}{(5)^2} = \frac{125}{25} = 5\text{ cm}
\end{align*}

Radius $r = 5\text{ cm}$ and Height $h = 5\text{ cm}$.

\item

Using integration, find the area lying above x-axis and included between the circle $x^{2}+y^{2}=8x$ and inside the parabola $y^{2}=4x$

\solution 

\begin{align*}
\text{Circle equation: } x^2 + y^2 &= 8x \implies (x-4)^2 + y^2 = 16 \quad \text{(Center: }(4,0), \text{Radius: } 4) \\
\text{Parabola equation: } y^2 &= 4x
\end{align*}

Substitute $y^2 = 4x$ into the circle equation to find intersection points:
\begin{align*}
x^2 + 4x &= 8x \\
x^2 - 4x &= 0 \implies x(x-4) = 0 \\
x &= 0 \text{ or } x = 4
\end{align*}
For $x = 0, y = 0$, and for $x = 4, y = \pm 4$.

The required area above the $x$-axis is split into two regions:
1. From $x = 0$ to $x = 4$, bounded by the parabola $y = 2\sqrt{x}$.
2. From $x = 4$ to $x = 8$, bounded by the circle $y = \sqrt{8x - x^2}$.

\begin{align*}
\text{Area} &= \int_{0}^{4} 2\sqrt{x}\,dx + \int_{4}^{8} \sqrt{8x - x^2}\,dx \\
&= 2 \left[ \frac{x^{3/2}}{3/2} \right]_{0}^{4} + \left[ \frac{x-4}{2}\sqrt{8x-x^2} + \frac{16}{2}\sin^{-1}\left(\frac{x-4}{4}\right) \right]_{4}^{8} \\
&= \frac{4}{3}(4^{3/2}) + \left[ 0 + 8\sin^{-1}(1) \right] - \left[ 0 + 8\sin^{-1}(0) \right] \\
&= \frac{32}{3} + 8\left(\frac{\pi}{2}\right) \\
&= \left(\frac{32}{3} + 4\pi\right) \text{ sq. units}
\end{align*}

\begin{center}
    or
\end{center}
\textbf{Question :}

Using the method of integration, find the area of the triangle ABC, coordinates of whose vertices are $A(2,0)$, $B(4,5)$ and $C(6,3)$.   

\solution 

\textbf{Equation of line $AB$ passing through $A(2,0)$ and $B(4,5)$:}
\begin{align*}
y - 0 &= \frac{5 - 0}{4 - 2}(x - 2) \\
y &= \frac{5}{2}(x - 2) \quad \implies y_1 = \frac{5}{2}(x - 2)
\end{align*}

\textbf{Equation of line $BC$ passing through $B(4,5)$ and $C(6,3)$:}
\begin{align*}
y - 5 &= \frac{3 - 5}{6 - 4}(x - 4) \\
y - 5 &= \frac{-2}{2}(x - 4) \\
y - 5 &= -(x - 4) \quad \implies y_2 = 9 - x
\end{align*}

\textbf{Equation of line $AC$ passing through $A(2,0)$ and $C(6,3)$:}
\begin{align*}
y - 0 &= \frac{3 - 0}{6 - 2}(x - 2) \\
y &= \frac{3}{4}(x - 2) \quad \implies y_3 = \frac{3}{4}(x - 2)
\end{align*}

\begin{align*}
\text{Area}(\triangle ABC) &= \text{Area under } AB \, (x \in [2,4]) + \text{Area under } BC \, (x \in [4,6]) - \text{Area under } AC \, (x \in [2,6]) \\
&= \int_{2}^{4} \frac{5}{2}(x - 2)\,dx + \int_{4}^{6} (9 - x)\,dx - \int_{2}^{6} \frac{3}{4}(x - 2)\,dx
\end{align*}

Evaluating the integrals individually:
\begin{align*}
I_1 &= \frac{5}{2} \left[ \frac{(x-2)^2}{2} \right]_{2}^{4} = \frac{5}{2} \left( \frac{4}{2} - 0 \right) = 5 \\
I_2 &= \left[ 9x - \frac{x^2}{2} \right]_{4}^{6} = \left( 54 - 18 \right) - \left( 36 - 8 \right) = 36 - 28 = 8 \\
I_3 &= \frac{3}{4} \left[ \frac{(x-2)^2}{2} \right]_{2}^{6} = \frac{3}{4} \left( \frac{16}{2} - 0 \right) = 6
\end{align*}

Combining the results:
\begin{align*}
\text{Area}(\triangle ABC) &= I_1 + I_2 - I_3 \\
&= 5 + 8 - 6 \\
&= 7 \text{ sq. units}
\end{align*}

area $\implies 7 \text{ sq. units}$

\item

If $A=[\begin{matrix}5&-1&4\\ 2&3&5\\ 5&-2&6\end{matrix}]$ find $A^{-1}$ and use it to solve the following system of equations:
$5x-y+4z=5$\newline
$2x+3y+5z=2$\newline 
$5x-2y+6z=-1$

\solution 

\begin{align*}
|A| &= 5(18 - (-10)) - (-1)(12 - 25) + 4(-4 - 15) \\
&= 5(28) + 1(-13) + 4(-19) \\
&= 140 - 13 - 76 = 51 \neq 0 \text{ (non-singular)}
\end{align*}

Cofactors of $A$:
\begin{align*}
C_{11} &= 28, \quad C_{12} = 13, \quad C_{13} = -19 \\
C_{21} &= -2, \quad C_{22} = 10, \quad C_{23} = 5 \\
C_{31} &= -17, \quad C_{32} = -17, \quad C_{33} = 17
\end{align*}

\begin{align*}
\text{Adj}(A) &= \begin{pmatrix} 28 & -2 & -17 \\ 13 & 10 & -17 \\ -19 & 5 & 17 \end{pmatrix} \\
A^{-1} &= \frac{1}{|A|} \text{Adj}(A) = \frac{1}{51} \begin{pmatrix} 28 & -2 & -17 \\ 13 & 10 & -17 \\ -19 & 5 & 17 \end{pmatrix}
\end{align*}

The matrix equation is $AX = B$, where $X = \begin{pmatrix} x \\ y \\ z \end{pmatrix}$ and $B = \begin{pmatrix} 5 \\ 2 \\ -1 \end{pmatrix}$.
\begin{align*}
X &= A^{-1}B = \frac{1}{51} \begin{pmatrix} 28(5) - 2(2) - 17(-1) \\ 13(5) + 10(2) - 17(-1) \\ -19(5) + 5(2) + 17(-1) \end{pmatrix} \\
&= \frac{1}{51} \begin{pmatrix} 140 - 4 + 17 \\ 65 + 20 + 17 \\ -95 + 10 - 17 \end{pmatrix} = \frac{1}{51} \begin{pmatrix} 153 \\ 102 \\ -102 \end{pmatrix} = \begin{pmatrix} 3 \\ 2 \\ -2 \end{pmatrix}
\end{align*}

\begin{center}
    or 
\end{center}.


\textbf{Question :} If $x, y, z$ are different and 
$$\begin{vmatrix} x & x^2 & 1+x^3 \\ y & y^2 & 1+y^3 \\ z & z^2 & 1+z^3 \end{vmatrix} = 0$$
then using properties of determinants show that $1+xyz=0$

\textbf{solution :}

Let $\Delta = \begin{vmatrix} x & x^2 & 1+x^3 \\ y & y^2 & 1+y^3 \\ z & z^2 & 1+z^3 \end{vmatrix}$.

Using the property of determinants that allows splitting a column containing a sum of two terms:
$$\Delta = \begin{vmatrix} x & x^2 & 1 \\ y & y^2 & 1 \\ z & z^2 & 1 \end{vmatrix} + \begin{vmatrix} x & x^2 & x^3 \\ y & y^2 & y^3 \\ z & z^2 & z^3 \end{vmatrix}$$

Taking out $x, y, z$ common from the $1^{\text{st}}, 2^{\text{nd}},$ and $3^{\text{rd}}$ rows respectively in the second determinant:
$$\Delta = \begin{vmatrix} x & x^2 & 1 \\ y & y^2 & 1 \\ z & z^2 & 1 \end{vmatrix} + xyz \begin{vmatrix} 1 & x & x^2 \\ 1 & y & y^2 \\ 1 & z & z^2 \end{vmatrix}$$

In the first determinant, interchange columns to match the second determinant:
\begin{align*}
\begin{vmatrix} x & x^2 & 1 \\ y & y^2 & 1 \\ z & z^2 & 1 \end{vmatrix} &= -\begin{vmatrix} x & 1 & x^2 \\ y & 1 & y^2 \\ z & 1 & z^2 \end{vmatrix} \quad (C_2 \leftrightarrow C_3) \\
&= (-1)^2 \begin{vmatrix} 1 & x & x^2 \\ 1 & y & y^2 \\ 1 & z & z^2 \end{vmatrix} \quad (C_1 \leftrightarrow C_2) \\
&= \begin{vmatrix} 1 & x & x^2 \\ 1 & y & y^2 \\ 1 & z & z^2 \end{vmatrix}
\end{align*}

$$\Delta = \begin{vmatrix} 1 & x & x^2 \\ 1 & y & y^2 \\ 1 & z & z^2 \end{vmatrix} + xyz \begin{vmatrix} 1 & x & x^2 \\ 1 & y & y^2 \\ 1 & z & z^2 \end{vmatrix}$$

$$\Delta = (1 + xyz) \begin{vmatrix} 1 & x & x^2 \\ 1 & y & y^2 \\ 1 & z & z^2 \end{vmatrix}$$

Now, apply row operations to evaluate the determinant:
\begin{align}
R_2 &\to R_2 - R_1 \\
R_3 &\to R_3 - R_1
\end{align}

$$\Delta = (1 + xyz) \begin{vmatrix} 1 & x & x^2 \\ 0 & y-x & y^2-x^2 \\ 0 & z-x & z^2-x^2 \end{vmatrix}$$

Take $(y-x)$ common from $R_2$ and $(z-x)$ common from $R_3$:
$$\Delta = (1 + xyz)(y - x)(z - x) \begin{vmatrix} 1 & x & x^2 \\ 0 & 1 & y+x \\ 0 & 1 & z+x \end{vmatrix}$$

\begin{align*}
\Delta &= (1 + xyz)(y - x)(z - x) [1 \cdot (z + x - (y + x))] \\
&= (1 + xyz)(y - x)(z - x)(z - y) \\
&= (1 + xyz)(x - y)(y - z)(z - x)
\end{align*}

We are given that $\Delta = 0$:
$$(1 + xyz)(x - y)(y - z)(z - x) = 0$$

Since $x, y, z$ are distinct, $(x - y) \neq 0$, $(y - z) \neq 0$, and $(z - x) \neq 0$.

Therefore:
$$1 + xyz = 0$$



\item

A card from a pack of 52 cards is lost. From the remaining cards of the pack, two cards are drawn randomly one-by-one without replacement and are found to be both kings. Find the probability of the lost card being a king.   

\solution 

\begin{align*}
\text{Let } E_1 &= \text{Lost card is a king} \implies P(E_1) = \frac{4}{52} = \frac{1}{13} \\
\text{Let } E_2 &= \text{Lost card is not a king} \implies P(E_2) = \frac{48}{52} = \frac{12}{13} \\
\text{Let } A &= \text{Two cards drawn from remaining 51 cards are both kings}
\end{align*}

Using conditional probabilities:
\begin{align*}
P(A | E_1) &= \frac{\text{3 kings left}}{\text{51 cards left}} = \frac{3_{C_2}}{51_{C_2}} = \frac{3 \times 2}{51 \times 50} = \frac{6}{2550} = \frac{1}{425} \\
P(A | E_2) &= \frac{\text{4 kings left}}{\text{51 cards left}} = \frac{4_{C_2}}{51_{C_2}} = \frac{4 \times 3}{51 \times 50} = \frac{12}{2550} = \frac{2}{425}
\end{align*}

Using Bayes' Theorem to find $P(E_1 | A)$:
\begin{align*}
P(E_1 | A) &= \frac{P(E_1)P(A | E_1)}{P(E_1)P(A | E_1) + P(E_2)P(A | E_2)} \\
&= \frac{\frac{1}{13} \times \frac{1}{425}}{\frac{1}{13} \times \frac{1}{425} + \frac{12}{13} \times \frac{2}{425}} \\
&= \frac{1}{1 + 24} = \frac{1}{25}
\end{align*}
\textbf{Answer:} $\frac{1}{25}$.






\end{enumerate}


